\subsection{\texorpdfstring{The Simplicity of $A_n$}{The Simplicity of An}}

Let $G$ be a group and let $\Omega$ be an infinite set.

\begin{exercise}[4.6.1]
	Prove that $A_n$ does not have a proper subgroup of index $< n$ for all $n \geq 5$.
\end{exercise}

\begin{sol}
	Suppose $H \leq A_n$ is a proper subgroup of index $k < n$. Then $A_n$ acts on the left cosets of $H$ by left multiplication, inducing a homomorphism $\phi : A_n \to S_k$. Observe that $n/2 > 0$ for all $n \geq 5$. This lets us conclude that $\ker\phi$ must be non-trivial, for otherwise $\phi$ would be injective and $|A_n| \leq |S_k|$. However, $k < n$ implies $k \leq n - 1$, or $k! \leq (n - 1)! < n!/2 = |A_n|$, which is a contradiction. Since $\ker\phi$ is a non-trivial normal subgroup of $A_n$, then $\ker\phi = A_n$ since $A_n$ is simple for all $n \geq 5$. It follows that $\phi$ is the trivial homomorphism, so that $H = A_n$, which is a contradiction. Therefore, $A_n$ does not have a proper subgroup of index $< n$ for all $n \geq 5$.
\end{sol}

\begin{exercise}[4.6.2]
	Find all normal subgroups of $S_n$ for all $n \geq 5$.
\end{exercise}

\begin{sol}
	Observe some normal subgroups of $S_n$ are $1, A_n$, and $S_n$ for all $n \geq 5$. To see that these are the only normal subgroups of $S_n$, let $N \nsub S_n$ be a normal subgroup, and consider $N \cap A_n$. We claim that $N \cap A_n \nsub A_n$.

	Let $\sigma \in N \cap A_n$. Then for any $\tau \in A_n$, we have that $\tau \sigma \tau\inv \in A_n$ since $A_n$ is a subgroup of $S_n$. Moreover, $\tau \sigma \tau\inv \in N$ since $N \nsub S_n$. It follows that $\tau \sigma \tau\inv \in N \cap A_n$, so that $N \cap A_n \nsub A_n$. Since $A_n$ is simple for all $n \geq 5$, then $N \cap A_n$ is either $1$ or $A_n$.
	\begin{enumerate}
		\item If $N \cap A_n = A_n$, then $A_n \subseteq N$. By maximality of $A_n$ in $S_n$, then $N = A_n$ or $N = S_n$.
		\item If $N \cap A_n = 1$, consider $NA_n \leq S_n$ because $N \nsub S_n$ and $A_n \nsub S_n$. It follows that $|NA_n| = |N||A_n|$. Because $NA_n \leq S_n$, then $|NA_n| = |N||A_n| \leq |S_n|$, so that $|N| \leq 2$. If $|N| = 2$, then $N = \gen\rho$ for some order 2 element $\rho$. However, $N \nsub S_n$ implies that $\tau\rho\tau\inv = \rho$ for all $\tau \in S_n$, or that $\rho \in Z(S_n)$, which is a contradiction since $Z(S_n) = 1$ for all $n \geq 3$. Therefore, $|N| = 1$, so that $N = 1$. \qh
	\end{enumerate}
\end{sol}

\begin{exercise}[4.6.3]
	Prove that $A_n$ is the only proper subgroup of index $< n$ in $S_n$ for all $n \geq 5$.
\end{exercise}

\begin{sol}
	Let $H \leq S_n$ be a proper subgroup where $|S_n : H| = k < n$. Then $S_n$ acts on the left cosets of $H$ by left multiplication, inducing a homomorphism $\phi : S_n \to S_k$. Observe that $n/2 > 0$ for all $n \geq 5$. This lets us conclude that $\ker\phi$ must be non-trivial, for otherwise $\phi$ would be injective and $|S_n| \leq |S_k|$. However, $k < n$ implies $k \leq n - 1$, or $k! \leq (n - 1)! < n! = |S_n|$, which is a contradiction. Since $\ker\phi$ is a non-trivial normal subgroup of $S_n$, then $\ker\phi = A_n$ or $\ker\phi = S_n$ since these are the only normal subgroups of $S_n$ for all $n \geq 5$. If $\ker\phi = S_n$, then $\phi$ is the trivial homomorphism, so that $H = S_n$, which is a contradiction. Therefore, $\ker\phi = A_n$, so that $A_n \subseteq H$. By maximality of $A_n$ in $S_n$, then $H = A_n$. Therefore, $A_n$ is the only proper subgroup of index $< n$ in $S_n$ for all $n \geq 5$.
\end{sol}

\begin{exercise}[4.6.4]
	Prove that $A_n$ is generated by the set of all 3-cycles for each $n \geq 3$.
\end{exercise}

\begin{sol}
	Let $S$ be the set of all 3-cycles in $A_n$. Observe that $\gen S \leq A_n$ since $S \subseteq A_n$. To show the other direction, let $\sigma \in A_n$. We show that two transpositions in $\sigma$ can be replaced by 3-cycles. There are two cases to discuss: Let $(a\ b), (c\ d)$ be two transpositions in $\sigma$.
	\begin{itemize}
		\item If $\set{a, b} \cap \set{c, d} = \emptyset$, then we can replace $(a\ b)(c\ d)$ by the product of 3-cycles $(a\ c\ d)(a\ d\ c)$.
		\item If $\set{a, b} \cap \set{c, d} = \set{x}$ for some $x$, then we can replace $(a\ b)(c\ d)$ by the 3-cycle $(a\ c\ d)$ if $x = a$ or $x = c$, and by the 3-cycle $(b\ c\ d)$ if $x = b$ or $x = d$.
	\end{itemize}
\end{sol}

\begin{exercise}[4.6.5]
	Prove that if there exists a chain of subgroups $G_1 \leq G_2 \leq \cdots \leq G$ such that
	\[G = \bigcup_{i = 1}^\infty G_i\]
	and each $G_i$ is simple, then $G$ is simple.
\end{exercise}

\begin{sol}
	Let $N \nsub G$. We first note that $N \cap G_i \nsub G_i$ for every $i$ as follows: for some $x \in N \cap G_i$ and $g \in G_i$, then $gxg\inv \in G_i$ since $G_i$ is a subgroup of $G$. Moreover, $gxg\inv \in N$ since $N \nsub G$. It follows that $gxg\inv \in N \cap G_i$, so that $N \cap G_i \nsub G_i$.

	By simplicity of each $G_i$, then $N \cap G_i$ is either $1$ or $G_i$. If $N \cap G_i = G_i$ for some $i$, then $G_i \subseteq N$. In particular, we have for some $G_j$ with $j > i$ that $G_i \subseteq G_j \subseteq N$, so that $N = G$. If $N \cap G_i = 1$ for all $i$, then $N$ is the union of the trivial subgroup, so that $N = 1$. Therefore, in either case, $N$ is either $1$ or $G$, so that $G$ is simple.
\end{sol}

\begin{exercise}[4.6.6]
	Let $D$ be the subgroup of $S_\Omega$ consisting of permutations which move only a finite number of elements of $\Omega$ (described in \exref{4.3.17}), and let $A$ be the set of all elements $\sigma \in D$ such that $\sigma$ acts as an even permutation on the (finite) set of points it moves. Prove that $A$ is an infinite simple group. (Show that every pair of elements of $D$ lie in a finite simple subgroup of $D$.)
\end{exercise}

\begin{sol}
	Note that for $\sigma, \tau \in A$, then $\sigma\tau\inv$ is an even permutation on the set of points moved by $\sigma$ and $\tau$, so that $A$ is a subgroup of $D$. Moreover, $A$ is normal in $D$ since for any $\sigma \in A$ and any $\tau \in D$, then $\tau\sigma\tau\inv$ is an even permutation on the set of points moved by $\sigma$ and $\tau$, so that $\tau\sigma\tau\inv \in A$. Therefore, $A \nsub D$. Moreover, $A$ is infinite since $\Omega$ is infinite, so that there are infinitely many even permutations on $\Omega$ that move only a finite number of elements of $\Omega$.

	Since $\Omega$ is infinite, for each $n \geq 5$, we can consider the subset $\Omega_n$ such that $|\Omega_n| = n$, and $\Omega_n \subseteq \Omega_{n + 1}$ for all $n$. Define
	\[B_n = \set{\sigma \in A \mid M(\sigma) \subseteq \Omega_n}\]
	Consider the bijection $\phi : B_n \to A_n$ defined by $\phi(\sigma) = \sigma|_{\Omega_n}$, where $\sigma|_{\Omega_n}$ is the restriction of $\sigma$ to $\Omega_n$. It is clear that this is a homomorphism, hence $B_n \cong A_n$. Moreover, $B_n$ is simple since $A_n$ is simple for all $n \geq 5$. Considering the chain of subgroups $B_5 \leq B_6 \leq \cdots$, we see that
	\[A = \bigcup_{n = 5}^\infty B_n.\]
	By \exref{4.6.5}, $A$ is simple.
\end{sol}

\begin{exercise}[4.6.7]
	Under the notation of the preceding exercise, prove that if $H \nsub S_\Omega$ and $H \neq 1$, then $A \leq H$, i.e., $A$ is the unique (non-trivial) minimal normal subgroup of $S_\Omega$.
\end{exercise}

\begin{sol}
	Let $\sigma \in H$. Since $\sigma$ moves only a finite number of elements of $\Omega$, there exists some $\Gamma \subseteq \Omega$ with $|\Gamma| \geq 5$ such that $\sigma(\omega) = \omega$ for all $\omega \in \Omega - \Gamma$. If $\sigma$ is an odd permutation, let $x, y \in \Omega - \Gamma$ be distinct (which exist since $\Omega$ is infinite) and replace $\sigma$ with $\sigma(x\ y) \in H$ and $\Gamma$ with $\Gamma \cup \set{x, y}$. Thus we may assume without loss of generality that $\sigma$ is an even permutation fixing every point outside $\Gamma$. 

	Let $B$ be the set of elements of $A$ which fix every point outside $\Gamma$. Since $B \cong A_{|\Gamma|}$ and $|\Gamma| \geq 5$, $B$ is simple. Since $H \cap B \trianglelefteq B$ is non-trivial (it contains $\sigma$), simplicity of $B$ gives $B \leq H$. In particular, $H$ contains a 3-cycle $(\alpha\ \beta\ \gamma)$. Now let $(a\ b\ c)$ be any 3-cycle in $S_\Omega$. Since there exists $\tau \in S_\Omega$ with $\tau(\alpha) = a$, $\tau(\beta) = b$, $\tau(\gamma) = c$, and $H \trianglelefteq S_\Omega$, we have $\tau(\alpha\ \beta\ \gamma)\tau^{-1} = (a\ b\ c) \in H$. Since every element of $A$ is a product of 3-cycles and $H$ is a subgroup, we conclude $A \leq H$.
\end{sol}

\begin{exercise}[4.6.8]
	Under the notation of the preceding two exercises, prove that $|D| = |A| = |\Omega|$. Deduce that
	\[\text{if }S_\Omega \cong S_\Delta, \text{ then } |\Omega| = |\Delta|.\]
	(Use the fact that $D$ is generated by transpositions. You may assume that countable unions and finite direct products of sets of cardinality $|\Omega|$ also have cardinality $|\Omega|$.)
\end{exercise}

\begin{sol}
	We first show that $|D| = |\Omega|$. Since $D$ is generated by transpositions, every element of $D$ is a finite product of transpositions. For each $n \geq 1$, the set of products of exactly $n$ transpositions has cardinality at most $|\Omega|^n = |\Omega|$, since each transposition is determined by a pair of elements of $\Omega$. Since
	\[D = \bigcup_{n=1}^\infty D_n\]
	where $D_n$ is the set of products of at most $n$ transpositions, and a countable union of sets of cardinality $|\Omega|$ has cardinality $|\Omega|$, we get $|D| \leq |\Omega|$. Since the map $\omega \mapsto (\omega\ \alpha)$ for a fixed $\alpha \in \Omega$ is an injection $\Omega \to D$, we conclude $|D| = |\Omega|$.

    We now show that $|A| = |\Omega|$. Since $A \leq D$, we have $|A| \leq |D| = |\Omega|$. Conversely, fix distinct $x, y \in \Omega$. Then the map $\omega \mapsto (\omega\ x\ y)$ for $\omega \in \Omega - \set{x, y}$ is an injection $\Omega - \set{x, y} \to A$, so $|\Omega| \leq |A|$. Therefore $|A| = |\Omega|$.

    Finally, suppose $S_\Omega \cong S_\Delta$ via some isomorphism $\phi$. By the preceding exercise, $A$ is the unique minimal normal subgroup of $S_\Omega$, and similarly $A'$ is the unique minimal normal subgroup of $S_\Delta$. Since $\phi$ maps normal subgroups to normal subgroups, $\phi(A) = A'$, so $|A| = |A'|$, and therefore $|\Omega| = |\Delta|$.
\end{sol}